% !TeX program = pdflatex
% Документация за курсов проект. Съвместима с pdfLaTeX.
\documentclass[12pt,a4paper]{article}

\usepackage[utf8]{inputenc}
\usepackage{cmap}
\usepackage[T2A]{fontenc}
\usepackage[english,bulgarian]{babel}
\selectlanguage{bulgarian}
\usepackage{geometry}
\geometry{margin=2.5cm}
\usepackage{microtype}
\usepackage{setspace}
\onehalfspacing
\setlength{\parindent}{0pt}
\setlength{\parskip}{6pt plus 2pt}
\usepackage[table]{xcolor}
\usepackage{graphicx}
\usepackage{float}
\usepackage{caption}
\usepackage{amsmath}
\usepackage{longtable,booktabs,array,multirow}
\usepackage{etoolbox}
\AtBeginEnvironment{longtable}{\small}
\setlength{\LTleft}{0pt}
\setlength{\LTright}{0pt}
\usepackage{enumitem}
\setlist{itemsep=2pt,topsep=4pt}
\usepackage{fancyvrb}
\usepackage{fvextra}
\DefineVerbatimEnvironment{verbatim}{Verbatim}{breaklines=true,breakanywhere=true,fontsize=\small}
\usepackage[hidelinks]{hyperref}

\providecommand{\tightlist}{\setlength{\itemsep}{0pt}\setlength{\parskip}{0pt}}
\providecommand{\pandocbounded}[1]{#1}
\providecommand{\passthrough}[1]{#1}
\sloppy
\hbadness=10000
\hfuzz=5pt

\definecolor{bestgreen}{RGB}{218,245,222}
\definecolor{softyellow}{RGB}{255,246,205}
\definecolor{headerblue}{RGB}{35,67,97}
\definecolor{lightblue}{RGB}{232,240,248}
\definecolor{lightgray}{RGB}{245,245,245}

\newcommand{\metric}[1]{\textbf{#1}}
\newcommand{\maybeincludegraphics}[2][]{%
  \IfFileExists{#2}{\includegraphics[#1]{#2}}{\fbox{\parbox{0.85\linewidth}{Липсва файлът с изображението: \texttt{\detokenize{#2}}}}}%
}

\begin{document}

\begin{titlepage}
\centering
\vspace*{1.2cm}
{\Large Софийски университет ``Св. Климент Охридски''\\[0.2cm]}
{\large Факултет по математика и информатика\\[1.2cm]}

{\Huge\bfseries Извличане на същности и връзки от научни резюмета\\[0.6cm]}

{\Large Документация за курсов проект\\[1.2cm]}

\begin{tabular}{r >{\raggedright\arraybackslash}p{10.5cm}}
Изготвил: & Явор Чамов \\
Факултетен номер: & 4MI3400787 \\
Предмет: & Извличане на знания от текст \\
Преподаватели: & проф. д-р Иван Койчев, проф. д-р Преслав Наков,\newline гл. ас. Димитър Димитров \\
Задачи: & Разпознаване на именувани същности;\newline Извличане на връзки \\
Основен резултат: & PubMedBERT за разпознаване на именувани същности и BioLinkBERT за извличане на връзки \\
\end{tabular}

\vspace{1.2cm}

\noindent \textbf{Декларация за липса на плагиатство}\\[0.2cm]
Декларирам, че настоящата курсова работа е моя собствена разработка, като всички използвани източници, идеи, илюстрации, данни и програмен код от други автори са коректно цитирани. Разбирам, че при установяване на плагиатство работата може да бъде оценена със ``Слаб''.

\vspace{1.0cm}

\noindent \textbf{Дата:} \today \hfill \textbf{Подпис:} \rule{4cm}{0.4pt}

\vfill
{\large \today}

\end{titlepage}

\tableofcontents
\newpage

\section{Въведение}

Настоящият документ описва цялостната работа по курсовия проект върху GutBrainIE 2026. Задачата е свързана с автоматично извличане на знания от научни заглавия и резюмета от PubMed, посветени на връзката между чревната микробиота и заболявания на нервната система или психичното здраве.

Работата е насочена към две задачи. Първата е T611, при която системата трябва да открива споменавания на биомедицински същности и да им задава етикет. Втората е T621, при която трябва да се откриват връзки между вече намерени текстови споменавания. Така проектът преминава от разпознаване на отделни обекти към изграждане на малки твърдения от вида ``субект -- отношение -- обект''.

Целта на проекта не е само да се получи висок числов резултат, а да се изгради възпроизводима верига за обработка: четене на данни, проверка на отместванията, обучение, оценяване, сравнение на различни подходи, изводи от грешките и кратък потребителски изглед за представяне на резултатите.

\section{Описание на задачите}

\subsection{T611: разпознаване на именувани същности}

При T611 входът е заглавие и резюме на научна статия. Изходът е списък от същности. Всяка същност съдържа място на срещане, начална и крайна позиция в текста, самия текстов откъс и етикет. За тази задача не се изисква връзка към външен речник или понятие.

\begin{verbatim}
{
  "PMID": {
    "entities": [
      {
        "start_idx": 75,
        "end_idx": 82,
        "location": "title",
        "text_span": "patients",
        "label": "human"
      }
    ]
  }
}
\end{verbatim}

Наличните етикети включват заболявания и находки, химични вещества, микробиом, бактерии, хора, животни, гени, храни, лекарства, хранителни добавки, анатомични места, биомедицински похвати и статистически похвати. Точните имена на етикетите се запазват в изходните файлове, защото са част от официалния формат.

\subsection{T621: извличане на връзки между споменавания}

При T621 входът отново е статия, но задачата вече е да се намери връзка между две текстови споменавания. Всяка връзка се представя чрез субект, сказуемо и обект. Изходът съдържа текстовия откъс и етикета на субекта, вида на връзката, както и текстовия откъс и етикета на обекта.

\begin{verbatim}
{
  "PMID": {
    "mention_level_relations": [
      {
        "subject_text_span": "intestinal microbiome",
        "subject_label": "microbiome",
        "predicate": "located in",
        "object_text_span": "patients",
        "object_label": "human"
      }
    ]
  }
}
\end{verbatim}

Тази задача е по-трудна от T611, защото грешките от разпознаването на същности се пренасят към извличането на връзки. Затова в проекта са разгледани два режима: оценяване със златни същности и крайно оценяване с предсказани същности.

\section{Набор от данни}

Данните са разделени на златен, сребърен, сребърен от 2025 г., бронзов и развоен набор. Златният набор е най-надежден и служи като основа за обучение и вътрешна проверка. Сребърните набори се използват като допълнителни примери. Бронзовият набор е най-голям, но е по-шумен и затова не е използван като основна проверка.

\begin{longtable}{p{3.1cm}rrrrr}
\caption{Обобщение на използваните набори от данни.}\\
\toprule
\textbf{Набор} & \textbf{Статии} & \textbf{Същности} & \textbf{Връзки} & \textbf{Същности/статия} & \textbf{Връзки/статия} \\
\midrule
\endfirsthead
\toprule
\textbf{Набор} & \textbf{Статии} & \textbf{Същности} & \textbf{Връзки} & \textbf{Същности/статия} & \textbf{Връзки/статия} \\
\midrule
\endhead
Златен & 639 & 20 530 & 7 892 & 32.13 & 12.35 \\
Сребърен & 811 & 26 134 & 10 487 & 32.22 & 12.93 \\
Сребърен 2025 & 499 & 15 275 & 7 420 & 30.61 & 14.87 \\
Бронзов & 2 972 & 89 987 & 28 230 & 30.28 & 9.50 \\
Развоен & 80 & 2 521 & 1 139 & 31.51 & 14.24 \\
\bottomrule
\end{longtable}

Златният набор съдържа 639 статии и 20 530 същности. Развойният набор съдържа 80 статии и 2 521 същности. Тези стойности показват, че всяка статия съдържа средно много на брой същности и връзки. Това прави задачата подходяща за обучени езикови подходи, но изисква точна обработка на позициите в текста.

\section{Неравномерност на етикетите}

Разпределението на етикетите не е равномерно. Най-честият етикет за същности е \texttt{DDF}, а при връзките най-често се срещат \texttt{influence}, \texttt{located in}, \texttt{target} и \texttt{is linked to}. Това е важна причина да се следят едновременно Micro-F1 и Macro-F1. Micro-F1 дава общата успеваемост, докато Macro-F1 показва дали системата се справя и с редките етикети.

\begin{figure}[H]
\centering
\maybeincludegraphics[width=0.95\linewidth]{entity_distribution.png}
\caption{Разпределение на етикетите за същности.}
\end{figure}

\begin{figure}[H]
\centering
\maybeincludegraphics[width=0.95\linewidth]{relation_distribution.png}
\caption{Разпределение на видовете връзки.}
\end{figure}

\begin{longtable}{p{2.5cm}p{2.7cm}rrp{2.8cm}rr}
\caption{Най-чести етикети и степен на неравномерност по набори.}\\
\toprule
\textbf{Набор} & \textbf{Най-честа същност} & \textbf{Дял} & \textbf{Съотн.} & \textbf{Най-честа връзка} & \textbf{Дял} & \textbf{Съотн.} \\
\midrule
\endfirsthead
\toprule
\textbf{Набор} & \textbf{Най-честа същност} & \textbf{Дял} & \textbf{Съотн.} & \textbf{Най-честа връзка} & \textbf{Дял} & \textbf{Съотн.} \\
\midrule
\endhead
Златен & \texttt{DDF} & 33.7\% & 25.49 & \texttt{influence} & 17.5\% & 197.57 \\
Сребърен & \texttt{DDF} & 32.5\% & 21.69 & \texttt{influence} & 19.1\% & 501.75 \\
Сребърен 2025 & \texttt{DDF} & 36.6\% & 34.47 & \texttt{influence} & 23.7\% & 219.88 \\
Бронзов & \texttt{DDF} & 38.4\% & 30.21 & \texttt{influence} & 22.2\% & 3129.50 \\
Развоен & \texttt{DDF} & 31.5\% & 22.66 & \texttt{located in} & 17.2\% & 49.00 \\
\bottomrule
\end{longtable}

При бронзовия набор неравномерността на връзките е най-силна: най-честата връзка е хиляди пъти по-честа от най-рядката. Затова бронзовите примери са подходящи само за внимателни допълнителни опити, но не и за надеждна проверка.

\subsection{Най-чести същности и връзки}

\begin{minipage}{0.48\linewidth}
\begin{longtable}{lr}
\caption{Най-чести етикети за същности.}\\
\toprule
\textbf{Етикет} & \textbf{Брой} \\
\midrule
\texttt{DDF} & 56 350 \\
\texttt{chemical} & 22 720 \\
\texttt{microbiome} & 14 618 \\
\texttt{human} & 11 051 \\
\texttt{anatomical location} & 10 088 \\
\texttt{bacteria} & 9 837 \\
\texttt{biomedical technique} & 7 395 \\
\texttt{animal} & 5 837 \\
\bottomrule
\end{longtable}
\end{minipage}
\hfill
\begin{minipage}{0.48\linewidth}
\begin{longtable}{lr}
\caption{Най-чести видове връзки.}\\
\toprule
\textbf{Връзка} & \textbf{Брой} \\
\midrule
\texttt{influence} & 11 568 \\
\texttt{located in} & 7 445 \\
\texttt{target} & 7 116 \\
\texttt{is linked to} & 6 373 \\
\texttt{affect} & 5 187 \\
\texttt{impact} & 3 018 \\
\texttt{is a} & 2 777 \\
\texttt{interact} & 2 238 \\
\bottomrule
\end{longtable}
\end{minipage}

\section{Подготовка и проверка на данните}

Данните са прочетени основно от файловете с разделител \texttt{|}. Всички PubMed номера се обработват като низове, за да не се получат разминавания между статии и бележки. За T611 позициите се проверяват отделно в заглавието и в резюмето. Ако \texttt{location} е \texttt{title}, тогава \texttt{start\_idx} и \texttt{end\_idx} се отнасят към заглавието; ако е \texttt{abstract}, те се отнасят към резюмето.

Тази проверка е важна, защото при разпознаването на същности правилният текстов откъс не е достатъчен. За вярно предсказание трябва да съвпаднат номерът на статията, мястото, началната позиция, крайната позиция и етикетът. При извличането на връзки се сравняват текстовият откъс и етикетът на субекта, видът връзка, както и текстовият откъс и етикетът на обекта.

\section{Архитектура на решението}

Системата е изградена като последователна верига от стъпки:

\begin{enumerate}
\item четене на статии, същности и връзки;
\item проверка на отместванията и премахване на точни повторения;
\item начален речников подход за T611;
\item обучени подходи за T611 чрез токенно разпознаване;
\item създаване на кандидат-двойки за T621;
\item начален правилов подход за T621;
\item обучен класификатор на двойки за T621;
\item оценяване, сравнение и запис на изхода в нужния вид;
\item потребителски изглед за преглед на статии, същности, връзки и грешки.
\end{enumerate}

За T611 основният силен подход е токенно разпознаване с предварително обучени биомедицински езикови мрежи. За T621 основният подход е класификатор на двойки: за всяка възможна двойка същности се създава пример, в текста се отбелязват субектът и обектът, а системата избира вид връзка или липса на връзка.

\section{Разделяне на данните и опити}

Основната стратегия е развойният набор да не се използва за обучение. Той се пази за крайна проверка. Златният набор се разделя вътрешно на 85\% обучение и 15\% проверка. След това се правят опити с добавяне на сребърен и сребърен 2025 набор. Бронзовият набор не е основен източник за обучение, защото е шумен и силно неравномерен.

Разгледаните главни групи опити са:

\begin{itemize}
\item начални речникови и правилови подходи;
\item GLiNER за разпознаване на същности;
\item BioBERT, SciBERT и PubMedBERT за токенно разпознаване;
\item класификатори на двойки за връзки с PubMedBERT и BioLinkBERT;
\item крайна верига, при която T621 използва същностите, предсказани от T611;
\item малки опити с езикови помощници за сравнение, без да са основното решение.
\end{itemize}

\section{Мерки за оценяване}

Използвани са точност, пълнота и F1. Точността показва каква част от предсказанията са верни. Пълнотата показва каква част от златните примери са намерени. F1 е средна хармонична стойност между тях.

\begin{align*}
P &= \frac{TP}{TP + FP}, \\
R &= \frac{TP}{TP + FN}, \\
F1 &= \frac{2PR}{P + R}.
\end{align*}

Micro-F1 обединява всички верни, грешни и пропуснати примери преди изчисляването. Macro-F1 първо изчислява стойност по етикет, а после усреднява. При този проект Macro-F1 е особено важен, защото показва дали редките същности и връзки не се пренебрегват.

\section{Резултати за T611}

Таблица \ref{tab:ner-results} показва резултатите за T611 в реда на проведените опити. Най-добрият ред е оцветен. Началният речников подход е оставен в таблицата, защото показва долната граница и доказва нуждата от обучени подходи.

\begin{longtable}{p{3.6cm}p{3.0cm}p{3.2cm}rr} \caption{Най-добри резултати за T611 -- разпознаване на същности.}\label{tab:ner-results}\\ \toprule \textbf{Опит} & \textbf{Модел} & \textbf{Обучение} & \textbf{Micro-F1} & \textbf{Macro-F1} \\ \midrule \endfirsthead \toprule \textbf{Опит} & \textbf{Модел} & \textbf{Обучение} & \textbf{Micro-F1} & \textbf{Macro-F1} \\ \midrule \endhead NER-PubMedBERT-GS25 & PubMedBERT & gold+silver+silver\_2025 & 0.8092 & 0.7472 \\ NER-PubMedBERT-GS & PubMedBERT & gold+silver & 0.8034 & 0.7311 \\ NER-SciBERT-GS25 & SciBERT & gold+silver+silver\_2025 & 0.7880 & 0.7198 \\ NER-SciBERT-GS & SciBERT & gold+silver & 0.7838 & 0.7158 \\ NER-BioBERT-GS & BioBERT & gold+silver & 0.7686 & 0.6975 \\ NER-PubMedBERT-G & PubMedBERT & gold & 0.7664 & 0.6270 \\ NER-GLiNER-GS & GLiNER & gold+silver & 0.6868 & 0.5828 \\ NER-Dictionary & речник & gold & 0.0013 & 0.0005 \\ \bottomrule \end{longtable}

Първите редове са обучени токенни класификатори с различни предварително обучени биомедицински модели. GLiNER и речниковият подход са оставени като сравнителни базови линии: GLiNER показва силата на span-based подход без класическо BIO декодиране, а речникът показва колко недостатъчно е само точното съвпадение с видени откъси.

\begin{figure}[H]
\centering
\maybeincludegraphics[width=0.98\linewidth]{t611_ner_dev_f1_by_experiment.png}
\caption{Сравнение на Micro-F1 и Macro-F1 за опитите по T611 върху развойния набор.}
\label{fig:t611-ner-dev-f1-by-experiment}
\end{figure}

Най-добър резултат за T611 постига \texttt{ner\_pubmedbert\_gold\_silver\_silver\_2025}. Той достига \metric{Micro-F1 = 0.8092} и \metric{Macro-F1 = 0.7472}. Това е най-добрият общ резултат и същевременно най-добрият баланс между честите и редките етикети. Добавянето на сребърните данни помага силно, а добавянето на сребърния набор от 2025 г. дава малко допълнително подобрение за PubMedBERT.

GLiNER е полезен като сравнителен подход, но в тези опити остава под токенните подходи. Речниковият подход е твърде строг, защото разчита на точно срещане на вече видяни откъси и не може надеждно да открива нови изрази или различни граници.

\section{Резултати за T621}

T621 е оценена в два режима. Първият използва златни същности от развойния набор. Той показва доколко добре се справя самият подход за връзки. Вторият режим използва същностите, предсказани от T611. Той е по-реалистичен, защото отразява цялата верига.
\newline
\newline
\begin{longtable}{p{3.1cm}p{2.6cm}p{2.2cm}p{4.2cm}rr}
\caption{Най-добри резултати за T621 -- извличане на връзки.}\label{tab:re-results}\\
\toprule
\textbf{Опит} & \textbf{RE модел} & \textbf{RE обучение} & \textbf{Същности / NER вход} & \textbf{Micro-F1} & \textbf{Macro-F1} \\
\midrule
\endfirsthead

\toprule
\textbf{Опит} & \textbf{RE модел} & \textbf{RE обучение} & \textbf{Същности / NER вход} & \textbf{Micro-F1} & \textbf{Macro-F1} \\
\midrule
\endhead

RE-BioLinkBERT-Gold
& BioLinkBERT 
& gold+silver 
& златни dev същности 
& 0.5257 
& 0.4978 \\

RE-BioLinkBERT-E2E-PubMedBERT-GS
& BioLinkBERT 
& gold+silver 
& PubMedBERT NER, обучен с gold+silver 
& 0.3842 
& 0.3117 \\

RE-BioLinkBERT-E2E-PubMedBERT-GS25
& BioLinkBERT
& gold+silver 
& PubMedBERT NER, обучен с gold+silver+silver 2025 
& 0.3810 
& 0.3125 \\

\bottomrule
\end{longtable}

Последните два реда използват един и същ модел за извличане на връзки, но различен вход от T611. Първият краен вариант използва същности от PubMedBERT, обучен върху gold+silver, а вторият използва същности от PubMedBERT, обучен върху gold+silver+silver 2025.

\begin{figure}[H]
\centering
\maybeincludegraphics[width=0.98\linewidth]{t621_re_dev_f1_by_experiment.png}
\caption{Сравнение на Micro-F1 и Macro-F1 за опитите по T621 върху развойния набор.}
\label{fig:t621-re-dev-f1-by-experiment}
\end{figure}

Най-добрият резултат за T621 при златни същности е \texttt{re\_pair\_classifier\_biolinkbert\_gold\_silver\_gold\_entities}. Той достига \metric{Micro-F1 = 0.5257} и \metric{Macro-F1 = 0.4978}. Този резултат показва, че класификаторът на двойки с BioLinkBERT е подходящ за биомедицински връзки, когато същностите са дадени правилно.

Най-добрият краен резултат е \texttt{re\_pair\_classifier\_biolinkbert\_gold\_silver\_pubmedbert\_gold\_silver\_entities}. Той достига \metric{Micro-F1 = 0.3842} и \metric{Macro-F1 = 0.3117}. Това е по-ниско от оценяването със златни същности, защото пропуски и грешни граници от T611 водят до пропуснати или грешни двойки при T621.

Много близък краен вариант е \texttt{re\_pair\_classifier\_biolinkbert\_gold\_silver\_pubmedbert\_gold\_silver\_silver\_2025\_entities}. Той има \metric{Micro-F1 = 0.3810} и \metric{Macro-F1 = 0.3125}. Макар общият резултат да е малко по-нисък, Macro-F1 е малко по-висок, което означава, че при част от по-редките връзки поведението може да е по-добро.

\section{Анализ на грешките}

Основните грешки при T611 са свързани с границите на същностите и с близки по смисъл етикети. Например химично вещество и лекарство могат да бъдат объркани, ако текстът не дава достатъчен знак за употребата на веществото. Също така понякога системата открива само част от израза, например \texttt{microbiome} вместо \texttt{intestinal microbiome}. Такива грешки се отчитат като грешни, защото границите не съвпадат точно.

При T621 грешките са повече поради две причини. Първо, броят на възможните двойки същности в една статия е голям. Второ, дори когато двете същности са намерени правилно, видът на връзката може да е труден за разграничаване. Най-чести обърквания се очакват между близки отношения като \texttt{impact} и \texttt{influence}, както и при посоката на връзката.

\begin{longtable}{lll r}
\caption{Най-чести тройки от тип субект--връзка--обект.}\\
\toprule
\textbf{Субект} & \textbf{Връзка} & \textbf{Обект} & \textbf{Брой} \\
\midrule
\texttt{microbiome} & \texttt{is linked to} & \texttt{DDF} & 6 373 \\
\texttt{chemical} & \texttt{influence} & \texttt{DDF} & 5 857 \\
\texttt{DDF} & \texttt{target} & \texttt{human} & 5 217 \\
\texttt{DDF} & \texttt{affect} & \texttt{DDF} & 5 187 \\
\texttt{DDF} & \texttt{is a} & \texttt{DDF} & 2 777 \\
\texttt{dietary supplement} & \texttt{influence} & \texttt{DDF} & 2 615 \\
\texttt{bacteria} & \texttt{influence} & \texttt{DDF} & 2 126 \\
\texttt{DDF} & \texttt{target} & \texttt{animal} & 1 899 \\
\texttt{chemical} & \texttt{located in} & \texttt{anatomical location} & 1 821 \\
\texttt{DDF} & \texttt{strike} & \texttt{anatomical location} & 1 387 \\
\bottomrule
\end{longtable}

Крайното извличане на връзки зависи силно от T611. Ако същност липсва, нито една връзка с нея не може да бъде намерена. Ако границата или етикетът са грешни, правилна връзка може да бъде оценена като грешна. Затова в доклада се разграничават ``извличане на връзки със златни същности'' и ``крайно извличане на връзки''.

\section{Потребителски изглед за представяне}

За представяне на проекта е предвиден лек потребителски изглед. Основната му цел е да показва работата на системата, а не да бъде завършен продукт. Подходящ избор е Streamlit, защото позволява бързо да се показват таблици, цветно отбелязани същности, връзки и мерки.

Минималният изглед съдържа:

\begin{itemize}
\item избор на статия по PubMed номер;
\item показване на заглавие и резюме;
\item цветно отбелязване на същностите от T611;
\item таблица с връзките от T621;
\item сравнение със златните бележки, когато са налични;
\item табло с Micro-F1, Macro-F1, точност и пълнота;
\item преглед на грешни и пропуснати примери.
\end{itemize}

Този изглед е полезен за защита на проекта, защото числата от таблиците се свързват с видими примери от реални статии. Така се показва не само кой подход е най-добър, но и как точно изглеждат правилните и грешните предсказания.

\section{Обсъждане}

Резултатите показват, че за T611 най-добре работи PubMedBERT, обучен върху златен, сребърен и сребърен 2025 набор. Причината е, че този подход използва предварително научени биомедицински езикови представяния и получава достатъчно примери за различните етикети.

За T621 най-силен е BioLinkBERT в ролята на класификатор на двойки. Това е очаквано, защото задачата изисква разбиране на връзката между две същности в биомедицински контекст. Разликата между резултата със златни същности и крайния резултат показва, че подобряването на T611 ще има пряк положителен ефект върху T621.

Сребърните данни са полезни, но не всяко добавяне води до подобрение при всяка настройка. При T611 сребърният 2025 набор помага на PubMedBERT, но при крайната T621 верига вариантът с най-силната самостоятелна T611 система не е най-добър по Micro-F1. Това показва, че най-добрият извличач на същности не винаги дава най-добрата входна основа за извличане на връзки.

\section{Ограничения}

Проектът има няколко ограничения. Първо, оценяването изисква точно съвпадение на граници и етикети, което наказва дори малки разлики в текстовия откъс. Второ, редките етикети имат малко примери и затова Macro-F1 остава по-нисък от Micro-F1. Трето, при T621 броят на кандидат-двойките е голям и се налага внимателно избиране на отрицателни примери. Четвърто, бронзовият набор е голям, но шумен; употребата му без контрол може да влоши резултатите.

\section{Заключение}

В проекта е изградена пълна верига за извличане на същности и връзки от биомедицински научни резюмета. Най-добрият подход за T611 е PubMedBERT, обучен върху златни и сребърни данни, включително сребърния набор от 2025 г. Най-добрият подход за T621 при златни същности е BioLinkBERT класификатор на двойки, обучен върху златен и сребърен набор. Най-добрият краен подход за T621 използва BioLinkBERT за връзки и същности, предсказани от PubMedBERT.

Основният извод е, че качеството на разпознаването на същности е решаващо за крайното извличане на връзки. Затова бъдеща работа трябва да се насочи към по-добра обработка на границите, повече примери за редките етикети, по-добро филтриране на кандидат-двойките и по-добър анализ на грешките по вид същност и вид връзка.

\section*{Източници}
\addcontentsline{toc}{section}{Източници и използвани файлове}

\begin{itemize}
\item Официална страница на GutBrainIE 2026: 
\url{https://hereditary.dei.unipd.it/challenges/gutbrainie/2026/}

\item Официално хранилище с начални решения: 
\url{https://github.com/MMartinelli-hub/GutBrainIE_2026_Baseline}

\item Yu Gu, Robert Tinn, Hao Cheng, Michael Lucas, Naoto Usuyama, Xiaodong Liu, Tristan Naumann, Jianfeng Gao, Hoifung Poon. 
\textit{Domain-Specific Language Model Pretraining for Biomedical Natural Language Processing}. 
ACM Transactions on Computing for Healthcare, 2021. 
Статията описва PubMedBERT/BiomedBERT, използван като най-добър модел за разпознаване на именувани същности в проекта. 
\url{https://arxiv.org/abs/2007.15779}

\item Michihiro Yasunaga, Jure Leskovec, Percy Liang. 
\textit{LinkBERT: Pretraining Language Models with Document Links}. 
ACL 2022. 
Статията описва LinkBERT и неговия биомедицински вариант BioLinkBERT, използван като най-добър модел за извличане на връзки в проекта. 
\url{https://aclanthology.org/2022.acl-long.551/}

\item Jinhyuk Lee, Wonjin Yoon, Sungdong Kim, Donghyeon Kim, Sunkyu Kim, Chan Ho So, Jaewoo Kang. 
\textit{BioBERT: a pre-trained biomedical language representation model for biomedical text mining}. 
Bioinformatics, 2020. 
Използван като сравнителен биомедицински модел за разпознаване на същности. 
\url{https://academic.oup.com/bioinformatics/article/36/4/1234/5566506}

\item Iz Beltagy, Kyle Lo, Arman Cohan. 
\textit{SciBERT: A Pretrained Language Model for Scientific Text}. 
EMNLP-IJCNLP 2019. 
Използван като сравнителен модел, предварително обучен върху научни текстове. 
\url{https://aclanthology.org/D19-1371/}

\item Jacob Devlin, Ming-Wei Chang, Kenton Lee, Kristina Toutanova. 
\textit{BERT: Pre-training of Deep Bidirectional Transformers for Language Understanding}. 
NAACL 2019. 
Основна архитектура, върху която са изградени използваните модели от семейството на BERT. 
\url{https://aclanthology.org/N19-1423/}
\end{itemize}
\end{document}
